\section{GoalScope}
\label{app:goalscope}
\subsection{Approach}
\label{app:goalscope:approach}

\subsubsection{Readout prompts}
\Cref{fig:readout-prompt} shows the complete prompts to probe the \emph{faithfulness} and \emph{causality} readouts throughout the paper. Neither probe receives the source prompt being probed as input. The faithfulness prompt begins with an empty assistant message and the causality prompt prefills the assistant response with a fixed C++ function stub.

\begin{figure}[ht]\centering
\includegraphics[width=0.85\linewidth]{Figures/goalscope/readout_prompts.png}
\caption{The full prompts used to measure the \emph{faithfulness} and \emph{causality} throughout this paper. We start our steering intervention at the last token position, and keep performing the same intervention throughout the generation.}
\label{fig:readout-prompt}
\end{figure}

% The following are the complete role-separated probe contents before applying the model family's chat template. Neither probe receives the source prompt as text. The faithfulness probe begins with an empty assistant message; the causality probe continues the fixed function prefix. The runner applies the template with \texttt{continue\_final\_message=True}.

% \paragraph{Faithfulness.}
% \PromptInput{faithfulness_probe.txt}

% \paragraph{Causality.}
% \PromptInput{causality_probe.txt}

\subsubsection{Intervention setup}

Following the notation used throughout the paper we call the source context being probed as the source prompt $p_\mathrm{src}$ and probe prompt as $p_\mathrm{probe}$. For a selected goal heads $\mathcal{G}$, we cache the output of each these heads at a specific token $\mathrm{tok}_{\mathrm{src}}$ being probed in the source run $M(p_\mathrm{src})$. We explore two intervention strategies using these head outputs.

\paragraph{additive\_norm.} We aggregate the head outputs into a single vector $v_g=\sum_{[\ell,j]\in H}a_t^{\ell j}$. We use  $v_g$ as a steering vector, adding it to a range of layers in the probe run $M(p_\mathrm{probe})$. Concretely, we perform the following intervention for each layer $\ell$ in the selected layer range.

\begin{equation}
h^{\ell}_{-1}\leftarrow h^{\ell}_{-1}+\alpha\|h^{\ell}_{-1}\|\frac{g_t}{\|g_t\|}.
\label{eqn:additive-norm}
\end{equation}
\arnab{does this eqn becomes recursive now? or is it okay?}

We perform the same injection repeatedly with greedy decoding. The steering vector direction $v_g$ remains fixed, with the coefficient $\alpha$ rescaling the magnitude at each intervention layer.

% and repeated injection during greedy decoding, with a 64-token generation budget. The source vector remains fixed. This coefficient specifies a fraction of the probe residual norm, rather than an absolute activation scale.

\paragraph{replace.} This intervention variant instead replaces the goal heads contributions at every generation step of the probe run $M(p_\mathrm{probe})$ with their corresponding contributions from the destination run. Concretely:

\begin{equation}
a^{\ell j}_{-1}\leftarrow \alpha \; a^{\ell j}_\mathrm{src}\;\text{, for each  } a^{\ell j} \in \mathcal{G}
\label{eqn:replace}
\end{equation}

We use $\alpha=1.0$ to denote exact replacement, with larger values denoting \textit{boosting} the head contributions to the residual stream at their corresponding places. 

We perform the same intervention (in \Cref{eqn:additive-norm} or \ref{eqn:replace}) at each generation step with greedy decoding. The probe prompt does not contain any information about the source context, so the information is communicated solely with the intervention. The main Qwen2.5-Coder-32B model configuration uses the \emph{additive\_norm} strategy with $\alpha=.1$, layers 35--45 inclusive. Some of the cross model comparisons use the \emph{replace} strategy with $\alpha=1.0$. We did not perform a systematic grid search to select these hyperparameter values.

% The replacement variant instead operates on each selected head's activation before its output projection, at that head's native layer. With a binary selection mask $m_h$, its update is $z_h\leftarrow(1-m_h)z_h+\alpha m_h z_h^{\mathrm{src}}$. We use $\alpha=1$ for the cross-model comparisons in \Cref{app:model_comparison}. Replacement does not use the additive layer band or normalize the source vector.

% \begin{center}\begin{minipage}{\linewidth}\centering\small
% \begin{tabular}{lp{.31\linewidth}p{.31\linewidth}}\toprule
%  & Norm-relative addition & Head replacement\\\midrule
% Cached signal & Sum of projected head contributions & Each selected head's pre-projection activation\\
% Injection sites & Probe layers 35--45 & Selected heads' own layers\\
% Source coefficient $\alpha$ & .1 & 1\\
% Duration & All decoding steps & All decoding steps\\
% \bottomrule\end{tabular}
% \captionof{table}{The two GoalScope intervention configurations. The additive settings are for the main Qwen2.5-Coder-32B evaluation; replacement is used for within-model replication across families and sizes.}
% \label{tab:steering-config}
% \end{minipage}\end{center}

% These settings were not selected by a systematic hyperparameter search. The localization sparsity coefficient $\lambda_{\mathrm{sparse}}=20$ (\Cref{app:head_loc}) is shared by the head sets used here; it is separate from the readout coefficient $\alpha$. The source and head-budget comparisons in \Cref{app:readout_sources} hold the additive configuration fixed. OVReader uses $\alpha=.15$ for transformed residual states and $.1$ for its head-output reference (\Cref{app:ovreader}).

\subsection{Evaluation Setup}

Our readout setup generates multiple tokens. We use an LLM judge (\texttt{claude-sonnet-4-6}) to check if the readout aligns with the LM's generated code for the frontier $g_k$. We include the most interesting bits of the faithfulness prompt in \Cref{fig:oracle-faith}. The causality prompt is almost identical, just replacing ``description'' with ``C++ code'' with some minor adaptions.   


\begin{figure}[ht]\centering
\includegraphics[width=\linewidth]{Figures/goalscope/evaluation/faithfulness.png}
\caption{Prompt for the LLM judge to evaluate \emph{faithfulness}. Note that we ask the judge to make a subjective decision of whether the readout is a non-trivial sub-part of the LM's generated code in response to $g_k$.}
\label{fig:oracle-faith}
\end{figure}

We count the judge response of \texttt{YES} as success and all other responses, including cut-off/corrupted responses from the judge, as failures. Note that in the oracle prompt we ask the judge to make a subjective judgement whether the readout captures a non-trivial portion of the upcoming LLM-generated code. So, the readout rates do not reflect $100\%$ alignment with the code (or execution accuracy of the C++ code in causality). 